
\subsection{Mencari Invers Suatu Matriks}

Dalam Contoh \ref{exa:verifyinginverse}, kita diberi $A^{-1}$ dan diminta memeriksa bahwa matriks ini
memang invers dari $A$. Pada bagian ini, kita membahas cara mencari $A^{-1}$. 

Misalkan
\begin{equation*}
A=\leftB
\begin{array}{rr}
1 & 1 \\
1 & 2
\end{array}
\rightB
\end{equation*}
seperti dalam Contoh \ref{exa:verifyinginverse}.
Untuk mencari $A^{-1}$, kita perlu mencari
matriks $\leftB
\begin{array}{rr}
x & z \\
y & w
\end{array}
\rightB $ sedemikian sehingga
\begin{equation*}
\leftB
\begin{array}{rr}
1 & 1 \\
1 & 2
\end{array}
\rightB \leftB
\begin{array}{rr}
x & z \\
y & w
\end{array}
\rightB =\leftB
\begin{array}{rr}
1 & 0 \\
0 & 1
\end{array}
\rightB 
\end{equation*}
Kita dapat mengalikan kedua matriks ini dan melihat bahwa agar persamaan tersebut benar, kita harus mencari solusi untuk sistem-sistem persamaan
\begin{equation*}
\begin{array}{c}
x+y=1 \\
x+2y=0
\end{array}
\end{equation*}
dan
\begin{equation*}
\begin{array}{c}
z+w=0 \\
z+2w=1
\end{array}
\end{equation*}
Menuliskan matriks diperbesar untuk kedua sistem ini menghasilkan
\begin{equation*}
\leftB
\begin{array}{rr|r}
1 & 1 & 1 \\
1 & 2 & 0
\end{array}
\rightB  
\label{inverse1a}
\end{equation*}
untuk sistem pertama dan
\begin{equation}
\leftB
\begin{array}{rr|r}
1 & 1 & 0 \\
1 & 2 & 1
\end{array}
\rightB  \label{inverse2a}
\end{equation}
untuk sistem kedua. 

Mari selesaikan sistem pertama. Kalikan baris pertama dengan $-1$,
lalu tambahkan hasilnya ke baris kedua untuk memperoleh
\begin{equation*}
\leftB
\begin{array}{rr|r}
1 & 1 & 1 \\
0 & 1 & -1
\end{array}
\rightB
\end{equation*}
Sekarang, kalikan baris kedua dengan $-1$, lalu tambahkan hasilnya ke baris pertama untuk
memperoleh
\begin{equation*}
\leftB
\begin{array}{rr|r}
1 & 0 & 2 \\
0 & 1 & -1
\end{array}
\rightB 
\end{equation*}
Jika ditulis berdasarkan variabel, ini menyatakan $x=2$ dan $y=-1.$

Sekarang, selesaikan sistem kedua, \ref{inverse2a}, untuk mencari $z$ dan $w.$ Anda akan memperoleh bahwa
$z = -1$ dan $w = 1$.

Jika nilai-nilai $x,y,z,$ dan $w$ yang diperoleh
dimasukkan ke dalam matriks invers, kita melihat bahwa inversnya
adalah
\begin{equation*}
A^{-1} = 
\leftB
\begin{array}{rr}
x & z \\
y & w
\end{array}
\rightB
=
\leftB
\begin{array}{rr}
2 & -1 \\
-1 & 1
\end{array}
\rightB 
\end{equation*}

Setelah meluangkan waktu untuk menyelesaikan sistem kedua, Anda mungkin memperhatikan bahwa operasi baris yang persis sama
digunakan untuk menyelesaikan kedua sistem. Dalam setiap kasus, hasil akhirnya berbentuk
$\leftB I|X\rightB $, dengan $I$ sebagai
identitas dan $X$ memberikan satu kolom invers. Pada pembahasan di atas,
\begin{equation*}
\leftB
\begin{array}{c}
x \\
y
\end{array}
\rightB
\end{equation*}
yang merupakan kolom pertama invers diperoleh dengan menyelesaikan sistem pertama, lalu
kolom kedua
\begin{equation*}
\leftB
\begin{array}{c}
z \\
w
\end{array}
\rightB 
\end{equation*}

Untuk menyederhanakan prosedur ini, kita sebenarnya dapat menyelesaikan kedua sistem sekaligus!
Untuk melakukannya, kita dapat menulis
\begin{equation*}
\leftB
\begin{array}{rr|rr}
1 & 1 & 1 & 0 \\
1 & 2 & 0 & 1
\end{array}
\rightB
\end{equation*}

dan melakukan reduksi baris hingga diperoleh
\begin{equation*}
\leftB
\begin{array}{rr|rr}
1 & 0 & 2 & -1 \\
0 & 1 & -1 & 1
\end{array}
\rightB
\end{equation*}
lalu membaca inversnya sebagai matriks $2\times 2$ di sisi kanan.

Pembahasan ini memotivasi algoritma penting berikut.

\begin{algorithm}{Algoritma Invers Matriks}{matrixinversionalgorithm}
Misalkan $A$ adalah matriks $n\times n$. \index{matriks!mencari invers} Untuk mencari $A^{-1}$ jika ada, bentuklah
matriks diperbesar $n\times 2n$
\begin{equation*}
\leftB A|I\rightB
\end{equation*}
Jika memungkinkan, lakukan operasi baris hingga diperoleh matriks $n\times 2n$
berbentuk
\begin{equation*}
\leftB I|B\rightB  
\end{equation*}
Jika hal ini berhasil dilakukan, $B=A^{-1}.$ Dalam kasus ini, kita mengatakan bahwa $A$ \textbf{invertibel}.
Jika reduksi baris menjadi
matriks berbentuk $\leftB I|B\rightB ,$ tidak mungkin dilakukan, maka $A$ tidak memiliki invers.
\end{algorithm}

Algoritma ini menunjukkan cara mencari invers jika invers tersebut ada. Algoritma ini juga akan memberi tahu Anda jika
$A$ tidak memiliki invers. 

Perhatikan contoh berikut.

\begin{example}{Mencari Invers}\label{exa:findinginverse}
Misalkan $A=\leftB
\begin{array}{rrr}
1 & 2 & 2 \\
1 & 0 & 2 \\
3 & 1 & -1
\end{array}
\rightB $. Carilah $A^{-1}$ jika ada.
\end{example}

\begin{solution} Bentuk matriks diperbesar
\begin{equation*}
\leftB A|I\rightB =  
\leftB
\begin{array}{rrr|rrr}
1 & 2 & 2 & 1 & 0 & 0 \\
1 & 0 & 2 & 0 & 1 & 0 \\
3 & 1 & -1 & 0 & 0 & 1
\end{array}
\rightB
\end{equation*}

Sekarang kita melakukan reduksi baris dengan tujuan memperoleh matriks identitas $3 \times 3$ di sisi kiri.
Pertama, kalikan baris pertama dengan $-1 $, lalu tambahkan hasilnya ke baris kedua;
setelah itu, kalikan baris pertama dengan $-3 $ dan tambahkan hasilnya ke baris ketiga. Proses ini
menghasilkan
\begin{equation*}
\allowbreak \leftB
\begin{array}{rrr|rrr}
1 & 2 & 2 & 1 & 0 & 0 \\
0 & -2 & 0 & -1 & 1 & 0 \\
0 & -5 & -7 & -3 & 0 & 1
\end{array}
\rightB 
\end{equation*}
Kemudian, kalikan baris kedua dengan 5 dan tambahkan hasilnya ke baris ketiga yang dikalikan dengan -2.
\begin{equation*}
\leftB
\begin{array}{rrr|rrr}
1 & 2 & 2 & 1 & 0 & 0 \\
0 & -10 & 0 & -5 & 5 & 0 \\
0 & 0 & 14 & 1 & 5 & -2
\end{array}
\rightB
\end{equation*}
Selanjutnya, tambahkan baris ketiga ke baris pertama yang dikalikan dengan $ -7$.
Ini menghasilkan
\begin{equation*}
\leftB 
\begin{array}{rrr|rrr}
-7 & -14 & 0 & -6 & 5 & -2 \\
0 & -10 & 0 & -5 & 5 & 0 \\
0 & 0 & 14 & 1 & 5 & -2
\end{array}
\rightB 
\end{equation*}
Sekarang, kalikan baris kedua dengan $ -\frac{7}{5}$, lalu tambahkan hasilnya ke baris pertama.
\begin{equation*}
 \leftB
\begin{array}{rrr|rrr}
-7 & 0 & 0 & 1 & -2 & -2 \\
0 & -10 & 0 & -5 & 5 & 0 \\
0 & 0 & 14 & 1 & 5 & -2
\end{array}
\rightB 
\end{equation*}
Terakhir, bagi baris pertama dengan -7, baris kedua dengan -10, dan baris ketiga
dengan 14 sehingga diperoleh
\begin{equation*}
\leftB 
\begin{array}{rrr|rrr}
1 & 0 & 0 & -\vspace{0.05in}\frac{1}{7} & \vspace{0.05in}\frac{2}{7} & \vspace{0.05in}\frac{2}{7} \\
0 & 1 & 0 & \vspace{0.05in}\frac{1}{2} & -\vspace{0.05in}\frac{1}{2} & 0
\\
0 & 0 & 1 & \vspace{0.05in}\frac{1}{14} & \vspace{0.05in}\frac{5}{14} & -
\vspace{0.05in}\frac{1}{7}
\end{array}
\rightB 
\end{equation*}
Perhatikan bahwa sisi kiri matriks ini kini merupakan matriks identitas $3 \times 3$, yaitu $I_3$.
Oleh karena itu, inversnya adalah matriks $3 \times 3$ di sisi kanan, yang diberikan oleh
\begin{equation*}
\leftB
\begin{array}{rrr}
-\vspace{0.05in}\frac{1}{7} & \vspace{0.05in}\frac{2}{7} & \vspace{0.05in}
\frac{2}{7} \\
\vspace{0.05in}\frac{1}{2} & -\vspace{0.05in}\frac{1}{2} & 0 \\
\vspace{0.05in}\frac{1}{14} & \vspace{0.05in}\frac{5}{14} & -\vspace{0.05in}
\frac{1}{7}
\end{array}
\rightB
\end{equation*}
\end{solution}

Melalui algoritma ini, mungkin saja Anda menemukan bahwa sisi kiri tidak dapat
direduksi baris menjadi matriks identitas. Perhatikan contoh situasi ini berikut. 

\begin{example}{Matriks yang Tidak Memiliki Invers}\label{exa:matrixnoinverse}
Misalkan $A=\leftB
\begin{array}{rrr}
1 & 2 & 2 \\
1 & 0 & 2 \\
2 & 2 & 4
\end{array}
\rightB $. Carilah $A^{-1}$ jika ada.
\end{example}

\begin{solution} Tuliskan matriks diperbesar $\leftB A|I\rightB $
\begin{equation*}
\leftB
\begin{array}{rrr|rrr}
1 & 2 & 2 & 1 & 0 & 0 \\
1 & 0 & 2 & 0 & 1 & 0 \\
2 & 2 & 4 & 0 & 0 & 1
\end{array}
\rightB
\end{equation*}
dan lanjutkan dengan operasi baris untuk mencoba memperoleh $\leftB
I|A^{-1}\rightB .$ Kalikan baris pertama dengan $-1$, lalu tambahkan hasilnya ke
baris kedua. Kemudian, kalikan baris pertama dengan $ -2 $ dan tambahkan hasilnya ke
baris ketiga.
\begin{equation*}
\leftB 
\begin{array}{rrr|rrr}
1 & 2 & 2 & 1 & 0 & 0 \\
0 & -2 & 0 & -1 & 1 & 0 \\
0 & -2 & 0 & -2 & 0 & 1
\end{array}
\rightB
\end{equation*}
Selanjutnya, tambahkan $-1 $ kali baris kedua ke baris ketiga.
\begin{equation*}
\leftB 
\begin{array}{rrr|rrr}
1 & 2 & 2 & 1 & 0 & 0 \\
0 & -2 & 0 & -1 & 1 & 0 \\
0 & 0 & 0 & -1 & -1 & 1
\end{array}
\rightB
\end{equation*}
Pada titik ini, Anda dapat melihat bahwa tidak ada cara untuk memperoleh $I$ di
sisi kiri matriks diperbesar ini. Jadi, algoritma ini tidak mungkin
diselesaikan, sehingga invers dari $A$ tidak
ada. Dalam hal ini, kita mengatakan bahwa $A$ tidak invertibel.
\end{solution}

Jika algoritma menghasilkan invers untuk matriks semula, jawaban selalu dapat diperiksa.
Untuk melakukannya, gunakan metode yang ditunjukkan dalam Contoh \ref{exa:verifyinginverse}. Periksa bahwa hasil kali $AA^{-1}$ dan $A^{-1}A$ keduanya sama dengan
matriks identitas. Dengan metode ini, Anda selalu dapat memastikan bahwa $A^{-1}$ telah dihitung dengan benar!

Salah satu kegunaan invers matriks adalah mencari solusi suatu sistem persamaan linear.
Ingat dari Definisi \ref{def:matrixform} bahwa kita dapat menulis suatu sistem persamaan dalam bentuk matriks,
yaitu berbentuk $AX=B$. Misalkan Anda telah mencari
invers matriks, yaitu $A^{-1}$. Anda kemudian dapat mengalikan kedua sisi
persamaan ini di sebelah kiri dengan $A^{-1}$ dan menyederhanakannya untuk memperoleh
\begin{equation*}
\begin{array}{c}
\left( A^{-1} \right) AX =A^{-1}B \\
\left(A^{-1}A\right) X = A^{-1}B \\
IX = A^{-1}B \\
X = A^{-1}B
\end{array}
\end{equation*}
Oleh karena itu, kita dapat mencari $X$, yaitu solusi sistem, dengan menghitung $X=A^{-1}B$.
Perhatikan bahwa setelah
$A^{-1}$ diperoleh, Anda dapat dengan mudah memperoleh solusi untuk ruas kanan yang berbeda
($B$ yang berbeda). Solusinya selalu $A^{-1}B$.
 
Kita akan membahas metode pencarian solusi sistem ini dalam contoh berikut. 

\begin{example}{Menggunakan Invers untuk Menyelesaikan Sistem Persamaan}\label{exa:inversetosolvesystem}
 Perhatikan sistem
persamaan berikut. Gunakan invers dari matriks yang sesuai untuk memperoleh solusi
sistem ini.
\begin{equation*}
\begin{array}{c}
x+z=1 \\
x-y+z=3 \\
x+y-z=2
\end{array}
\end{equation*}
\end{example}

\begin{solution} Pertama, kita dapat menuliskan sistem persamaan dalam bentuk matriks
\begin{equation}
AX = 
\leftB
\begin{array}{rrr}
1 & 0 & 1 \\
1 & -1 & 1 \\
1 & 1 & -1
\end{array}
\rightB \leftB
\begin{array}{r}
x \\
y \\
z
\end{array}
\rightB =\leftB
\begin{array}{r}
1 \\
3 \\
2
\end{array}
\rightB  = B  \label{inversesystem1}
\end{equation}

Invers dari matriks
\begin{equation*}
A = \leftB
\begin{array}{rrr}
1 & 0 & 1 \\
1 & -1 & 1 \\
1 & 1 & -1
\end{array}
\rightB
\end{equation*}
adalah
\begin{equation*}
A^{-1} = \leftB
\begin{array}{rrr}
0 & \vspace{0.05in}\frac{1}{2} & \vspace{0.05in}\frac{1}{2} \\
1 & -1 & 0 \\
1 & -\vspace{0.05in}\frac{1}{2} & -\vspace{0.05in}\frac{1}{2}
\end{array}
\rightB
\end{equation*}

Pemeriksaan invers ini diserahkan sebagai latihan.

Dari sini, solusi sistem \ref{inversesystem1} yang diberikan
diperoleh melalui
\begin{equation*}
\leftB
\begin{array}{r}
x \\
y \\
z
\end{array}
\rightB = A^{-1}B = \leftB
\begin{array}{rrr}
0 & \vspace{0.05in}\frac{1}{2} & \vspace{0.05in}\frac{1}{2} \\
1 & -1 & 0 \\
1 & -\vspace{0.05in}\frac{1}{2} & -\vspace{0.05in}\frac{1}{2}
\end{array}
\rightB \leftB
\begin{array}{r}
1 \\
3 \\
2
\end{array}
\rightB =\leftB
\begin{array}{r}
\vspace{0.05in}\frac{5}{2} \\
-2 \\
-\vspace{0.05in}\frac{3}{2}
\end{array}
\rightB 
\end{equation*}
\end{solution}

Bagaimana jika ruas kanan, yaitu $B$, dari \ref{inversesystem1} adalah $\leftB
\begin{array}{r}
0 \\
1 \\
3
\end{array}
\rightB ?$ Dengan kata lain, apa solusi dari
\begin{equation*}
\leftB
\begin{array}{rrr}
1 & 0 & 1 \\
1 & -1 & 1 \\
1 & 1 & -1
\end{array}
\rightB \leftB
\begin{array}{r}
x \\
y \\
z
\end{array}
\rightB =\leftB
\begin{array}{r}
0 \\
1 \\
3
\end{array}
\rightB ?
\end{equation*}
Berdasarkan pembahasan di atas, solusinya diberikan oleh
\begin{equation*}
\leftB
\begin{array}{r}
x \\
y \\
z
\end{array}
\rightB = A^{-1}B = \leftB
\begin{array}{rrr}
0 & \vspace{0.05in}\frac{1}{2} & \vspace{0.05in}\frac{1}{2} \\
1 & -1 & 0 \\
1 & -\vspace{0.05in}\frac{1}{2} & -\vspace{0.05in}\frac{1}{2}
\end{array}
\rightB \leftB
\begin{array}{r}
0 \\
1 \\
3
\end{array}
\rightB =\leftB
\begin{array}{r}
2 \\
-1 \\
-2
\end{array}
\rightB 
\end{equation*}
Ini menunjukkan bahwa untuk sistem $AX=B$ dengan $A^{-1}$ yang ada,
solusi mudah dicari ketika vektor $B$ diubah.
%
%We conclude this section with some important properties of the inverse. 
%
\begin{theorem}{Invers dari Transpos dan Hasil Kali}\label{thm:inversetransposeprod}
%Let $A, B$, and $A_i$ for $i=1,...,k$ be $n \times n$ matrices. 
\begin{enumerate}
\item
%If $A$ is an invertible matrix, then $(A^{T})^{-1} = (A^{-1})^{T}$
\item
%If $A$ and $B$ are invertible matrices, then $AB$ is invertible and $(AB)^{-1} = B^{-1}A^{-1}$
\item
%If $A_1, A_2, ..., A_k$ are invertible, then the product $A_1A_2 \cdots A_k$ is invertible, and $(A_1A_2 \cdots A_k)^{-1} = A_k^{-1}A_{k-1}^{-1} \cdots A_2^{-1}A_1^{-1}$
\end{enumerate}
\end{theorem}
%
%Consider the following theorem.
%
\begin{theorem}{Sifat-Sifat Invers}\label{thm:inverseproperties}
%Let $A$ be an $n \times n$ matrix and $I$ the usual identity matrix. 
\begin{enumerate}
\item
%$I$ is invertible and $I^{-1} = I$
\item
%If $A$ is invertible then so is $A^{-1}$, and $(A^{-1})^{-1} = A$
\item
%If $A$ is invertible then so is $A^k$, and $(A^k)^{-1} = (A^{-1})^k$
\item
%If $A$ is invertible and $p$ is a nonzero real number, then $pA$ is invertible and $(pA)^{-1} = \frac{1}{p}A^{-1}$
\end{enumerate}
\end{theorem}
